\documentclass{article} % For LaTeX2e
\usepackage{iclr2027_conference,times}

\usepackage[T1]{fontenc}

% Optional math commands from https://github.com/goodfeli/dlbook_notation.
% This also pulls in amsmath, amsfonts and bm (see math_commands.tex line 3),
% which we rely on below for \text{...} inside inline math.
\input{math_commands.tex}

\usepackage{hyperref}
\usepackage{url}

\title{Confidence is not competence: inverse-folding models know binding in the mixed derivative}

% ANONYMOUS SUBMISSION. Do not add real names/emails/affiliations here.
% The iclr2027_conference.sty style automatically prints "Anonymous authors /
% Paper under double-blind review" in place of whatever \author{} contains,
% as long as \iclrfinalcopy (below) stays commented out. \author{} content is
% required by \maketitle but is not rendered in this (default) anonymous mode.
\author{Anonymous authors \\ Anonymous affiliation}

%\iclrfinalcopy % Uncomment for camera-ready version, but NOT for submission.

\begin{document}

\maketitle

\begin{abstract}
Staged binder design (generate a backbone, then inverse-fold a sequence) is thought to stumble at interface
\emph{hotspots}. A prominent report (ProBID-Net) quantifies this as inverse-folding recovery of 0.334 at hotspots
versus 0.472 elsewhere. The phenomenon is real but misread. \textbf{Confidence and competence are two different
derivatives of the same inverse-folding likelihood.} A residue's \emph{confidence} is the \emph{diagonal}, any
scalar of the bound-conditioned distribution $P$, and it measures fold-stability constraint. Its binding
\emph{leverage} is the \emph{mixed second derivative}, the response to ablating the partner. We prove $L$ is not a
function of $P$, so every scalar the field reads off these models is blind to binding \textbf{by construction}.
Empirically, a flexible learner over the entire bound distribution recovers only about 37\% of $L$. The
consequence is large and controlled. On natural complexes (SKEMPI) confidence sits at chance and adds nothing
beyond geometry, while the mixed derivative, a \textbf{zero-shot} readout that uses no binding labels, adds binding
information \textbf{beyond geometry, conservation, and the one-pass log-odds} (the full published feature set, and
it beats a \emph{supervised} baseline), across four inverse-folding architectures. The published deficit is
largely a burial confound. And the direction is \textbf{actionable}. An $+\alpha \cdot L$ tilt on a \emph{frozen}
ProteinMPNN steers it toward higher binding-favorability, confirmed against a matched random control by
independent sequence models, structure predictors, \textbf{and a physics energy function} (FoldX). The physics
function also shows the tilt beats a same-magnitude \emph{confidence} direction, isolating a binding-specific
effect (four judges, two folders, both steering directions, 120 folded complexes). The effect holds when the tilt
is applied on \emph{predicted}, not crystal, backbones, the staged-design regime itself. The leverage operator is
BA-Cycle (Jiao et al., 2024). Our contribution is the decomposition, the identifiability no-go, the first
beyond-geometry-and-conservation control on an inverse-folding binding signal, the feature-class law, and
frozen-model steering. $L$ is the model's \textbf{classifier-free-guidance direction}: to read an un-trained
quantity off a conditional generative model, ablate the conditioner and take the mixed derivative, not the
marginal mistaken for competence.
\end{abstract}

\section{Introduction}

The dominant paradigm for designing protein binders is staged: a generative model proposes a backbone, and an
inverse-folding model assigns a sequence by maximising $p(\text{sequence} \mid \text{backbone})$. This
factorisation is convenient, but it omits the very quantity a binder exists to optimise. There is no
binding-energy term anywhere in $p(\text{sequence} \mid \text{backbone})$. The concern this raises is sharpest at
\emph{interface hotspots}: the small set of residues that contribute most of the binding free energy, and which
are frequently \emph{frustrated}, buried polar residues or strained rotamers that are locally unfavourable for
the monomer's fold but bought because they pay off in binding. If inverse folding optimises fold-compatibility,
the reasoning goes, it should systematically miss exactly these residues, and a hotspot should sit in the tail of
the model's distribution rather than at its mode. A prominent measurement appears to confirm this: ProBID-Net
reports inverse-folding recovery of 0.334 at hotspots versus 0.472 elsewhere, and attributes the gap to dynamics.

This paper asks one question: when an inverse-folding model is used to build a binder, where, if anywhere, does
its knowledge of \emph{binding} live, and can a designer read it and act on it? The scope is deliberately narrow
and practical: structure-conditioned inverse-folding models (ProteinMPNN, ESM-IF1, and three others),
protein-protein interface binding, and readouts computable at design time. Our answer is that the binding signal
is real, but the field has been reading it from the wrong place. It has read \emph{scalar} summaries of the
model's output (recovery, confidence, a complex-vs-monomer KL) that are blind to binding \textbf{by
construction}, when it lives in a \textbf{mixed second derivative} that is cheap to compute and directly
actionable.

The consequence is immediate for practice, in three forms. As a \textbf{diagnostic}: inverse-folding confidence
is a geometry detector, not a binding score, so a pipeline that distrusts it at the interface (as the field's
leading one-shot designer does, by freezing the interface) is right to, and can now do better. As a
\textbf{training-free ranker}: the mixed derivative locates interface hotspots and adds on top of the exact
feature set published predictors already use, at no training cost. As a \textbf{drop-in steering knob}: a
$+\alpha \cdot L$ tilt biases a \emph{frozen, off-the-shelf} inverse-folding model toward higher-binding
interface residues with no retraining, confirmed by an independent sequence model, an independent structure
predictor, and a physics energy function. The recipe is not protein-specific. The mixed derivative is the model's
\textbf{classifier-free-guidance direction}, so the same move (ablate the conditioner, read the mixed derivative,
not the marginal) applies to any conditional generative model, which is what places the result at a
representation-learning venue rather than a purely methodological one.

We are precise about what is new, because the nearest prior methods already use this operator. The leverage
operator itself, a partner-ablated second forward pass, is the BA-Cycle operator of Jiao et al. (2024), and
decoders such as RedNet already tilt along this direction. Our contribution is not the operator but what it means
and what it buys. We give a decomposition that separates confidence from leverage at the mechanism level, a proof
that no scalar of the bound distribution can recover leverage, the first control showing the signal survives
beyond geometry and beyond evolutionary conservation (the full feature set published hotspot predictors use), a
feature-class law that holds across four inverse-folding architectures, and steering of a frozen model that three
independent instruments confirm. Where prior work uses the direction, we characterise it, bound what it can and
cannot express, and show that a designer can read it and act on it.

A single principle runs through all of this: confidence is not competence, and it recurs at every level of the
modeling stack. An inverse-folding model's confidence (any scalar of its bound distribution) is blind to binding
by construction. A structure predictor's interface ipTM, confidence one level up, is blind to
binding-\emph{specificity}. It rewards a confidently-packed interface, which a mere foldability tilt achieves, so
a be-more-confident tilt beats the binding direction on ipTM while losing to it under models with an explicit
binding term. It takes a readout with such a term, the mixed derivative or a physics energy function, to see
binding at all. The competence always lives in the response to ablating the partner, never in the bound-state
confidence, whichever model does the reading.

There is a confound that no prior analysis controls, and it runs opposite to intuition. Hotspots are, on average,
\emph{more deeply buried} than other interface residues, and burial is precisely where inverse folding is
\emph{most} confident and accurate. An uncontrolled hotspot-vs-rest comparison therefore mixes a putative binding
effect with a large burial effect of the opposite sign, so the naive comparison \textbf{hides} any real deficit
rather than inventing one. Controlling burial is not a detail. It is the experiment.

\paragraph{Contributions.}
\begin{itemize}
\item \textbf{A decomposition and a no-go theorem.} Confidence is the diagonal term of the inverse-folding
  log-likelihood and leverage is the mixed second derivative, and we prove leverage is not a function of the bound
  distribution. A flexible learner over the entire bound distribution recovers only about 37\% of it, so about
  63\% is irreducible from $P$ in both inverse-folding families.
\item \textbf{A feature-class law.} On natural complexes (SKEMPI) every scalar of the bound distribution
  (confidence, negentropy, the scalar KL) sits at or below a calibrated placebo floor, while the mixed derivative
  adds binding information beyond geometry, beyond evolutionary conservation, and beyond the one-pass log-odds,
  across four architectures.
\item \textbf{The published deficit is largely a burial confound.} Under a pre-registered matched-pair design the
  crystal-backbone hotspot deficit attenuates sharply across five architectures and under ProBID-Net's own
  released model.
\item \textbf{Steering a frozen model.} A $+\alpha \cdot L$ tilt on a frozen ProteinMPNN raises
  binding-favorability, confirmed by four inverse-folding judges, two structure predictors, and a physics energy
  function (FoldX), on crystal and on predicted backbones.
\item \textbf{Generality.} The blindness extends from binding to catalytic residues, and the recipe is the
  classifier-free-guidance direction for any conditional generative model.
\end{itemize}

\section{Setup and pre-registration}
% TODO (next assembly pass): migrate from PAPER_DRAFT.md \S2, house-styled.
% Fold: state complex-level, burial-matched leakage control plainly here (recon action).
% Keep the four evaluation questions (Q1-Q4) stated up front (recon: eval questions early).

\section{Confidence is not competence}
% TODO (next assembly pass): migrate from PAPER_DRAFT.md \S3, house-styled.

\section{The Confidence--Leverage Decomposition: the model knows binding in the mixed derivative}
% TODO (next assembly pass): migrate from PAPER_DRAFT.md \S4, house-styled.
% Fold: foreground the multi-oracle anti-circular evaluation (AF2/Boltz-2 ipTM + FoldX) earlier;
% add a methods line asserting identical baseline settings (temperature/decoding-order/sampling)
% for every comparator (recon actions).

\section{On crystal backbones, the hotspot gap is a burial artifact}
% TODO (next assembly pass): migrate from PAPER_DRAFT.md \S5, house-styled.

\section{The tax lives in the conditioning set}
% TODO (next assembly pass): migrate from PAPER_DRAFT.md \S6, house-styled.

\section{Ruling out competing mechanisms}
% TODO (next assembly pass): migrate from PAPER_DRAFT.md \S7, house-styled.

\section{Related work and positioning}
% TODO (next assembly pass): migrate from PAPER_DRAFT.md \S8, house-styled.
% Fold: put the mechanism-level "why not just BA-Cycle/RedNet" differentiation at the top (recon action).

\section{Limitations}
% TODO (next assembly pass): migrate from PAPER_DRAFT.md \S9, house-styled.
% Fold: add a one-paragraph "why the interim in-silico signal is informative" justification (recon action).

\section*{Reproducibility statement}
% TODO (next assembly pass): migrate from PAPER_DRAFT.md "Reproducibility and LLM-usage disclosure".

\section*{LLM usage disclosure}
% TODO (next assembly pass): migrate from PAPER_DRAFT.md "Reproducibility and LLM-usage disclosure".

\appendix

\section{Pre-registered false-positive modes and their controls}
% TODO (next assembly pass): migrate from PAPER_DRAFT.md "Appendix A".

\end{document}
